
\section{Non-detection of Low-Frequency Radio Emission from MTP0015 and MTP044} \label{sec:mtp}
Traditional radio pulsar search methods and survey specifications select against long-period and longer-duration single pulses~\citep{morello}. The MeerTRAP project~\citep{sanidas} runs a commensal real-time search for single-pulse emission on the MeerKAT telescope~\citep{jonas} which is well suited to find intermittent and long-period sources. MeerTRAP recently reported $26$ new discoveries~\citep{meertrap}. Two of these were North of $28\deg$ declination and appropriate for observation with the Irish LOFAR station which routinely observes such transient radio neutron stars~\citep{david}.

\textit{MTP0015}: this $1.077$-s source is seen with a pulse amplitude distribution suggestive of nulling with MeerKAT and has also been repeatedly detected with CHIME~\citep{dong2023}. In CHIME observations the spectrum appears brightest at the low-end of the band ($\sim 400-500$~MHz). 

%has been reguarly observed (paper) in radio since it was first detected on X (paper). It pulses 1.2 times/hr according to CHIME and MeerTRAP. The spin period has been characterized as 1.0xxx seconds with a reported dispersion measure (DM) of 38.1 \dmu. Observed pulses have shown brightness and microstructure with evidence of scintillation across the band. This makes MTP0015 an interesting target for low frequency observations. 

\textit{MTP0044}: this $17.495$-s source is so-far the third slowest spinning radio-emitting neutron star. It has so-far been seen exclusively in the MeerKAT UHF band ($544-1088$~MHz). 
%is a newly discovered source by MeerTRAP (paper). The source was detected 9 times during 1.8 hours of observations in UHF (544 -- 1088 MHz). MTP0044 is thought to be a repeating transient, further investigation of pulse profile and period to determine its specific class. Paper reports an estimated DM and period $55.8 \pm 0.04$ \dmu and $17.49487 \pm $ s. 

% \section{Observations}
Both sources were observed for $3$ hours, between 2025-01-15 and 2025-01-17, using the Irish LOFAR station observing in the $102.246338-197.558594$~MHz frequency range. This adds to previous observations (same specifications) of MTP0015 between July and December 2021, totalling $5$ hours~\citep{david_thesis}. The data were collected as described in \citet{david}; in particular coherent dedispersion was performed at the reported dispersion measure values and high-resolution Stokes $I$ filterbanks were created for post-processing.

A total of 4 hours of observations were carried out with I-LOFAR. Each target was observed for 2 hours (see \cref{tab:obs_details} for details)  each using the High Band Antenna (HBA; 110 -- 190 MHz). 

% \begin{table}[h]
%    \centering
%    \begin{tabular}{cccccc}
%        \hline 
%         Source  & RA (J2000) & DEC (J2000) & Period, $P$ & DM & Obs. MJDs \\
%         & (hh:mm:ss) & (dd:mm:ss) & (s) & (\dmn)& \\ \hline 
%         MTP0015 & 22:37:29.4 & +28:28:40 & & & \\
%         MTP0044 & 22:18:23.3 & +29:02:56 & & & \\ \hline \hline
%    \end{tabular}
%    \caption{Observational Parameters of MTP0015 and MTP0044}
%    \label{tab:obs_details}
% \end{table}

Data were cleaned for radio frequency interference (RFI) using \textsc{presto}~\citep{presto}. Single pulse searches were performed using \textsc{presto} but also \textsc{transientX}~\citep{transientx}, the latter of which performed additional RFI cleaning. Periodicity searching was performed by folding the data using the reported periodicities and, as there was no reference period epoch, scanning a range around that nominal value. No detections were made. The $8$-$
\sigma$ single pulse peak flux density limits are for MTP0014: 28.5, 24.1, 21.4, 19.8 $\times (W/10\;\mathrm{ms})^{-1/2}$~Jy; and for MTP0044: 
30.1, 24.7, 22.2, 20.3 $\times (W/10\;\mathrm{ms})^{-1/2}$~Jy; 
at 120, 140, 160 and 180 MHz respectively, where $W$ is the pulse width. The $8$-$
\sigma$ periodic average flux density limits are for MTP0014: 3.4, 2.9, 2.5, 2.4 mJy; for MTP0044: 3.6, 3.0, 2.6, 2.4 mJy; at 120, 140, 160 and 180 MHz respectively (assuming a 1\% duty cycle).